\section{Introduction}
\label{sec:intro}

\subsection{The Bloxorz lineage}
\label{sec:lineage}

Bloxorz is a flash puzzle game released in 2007 by Damien Clarke. The player rolls a
1$\times$1$\times$2 block over a 2D grid, dropping it upright into a target hole. Despite
the trivial action set (\emph{up}, \emph{down}, \emph{left}, \emph{right}), Bloxorz's
mechanics admit dramatically non-trivial reasoning: the block's orientation determines
which tiles it occupies, which switches it activates, and whether it falls through weak
tiles. Optimal solutions on later levels exceed 100 moves and require navigating
irreversible side-effects of switch activations.

Our levels are sourced from grahambarrgraham's open-source Bloxorz repository
\cite{grahambarrgrahambloxorz}, and our verifier is a re-implementation of tkoz0's
breadth-first search solver \cite{tkoz0bloxorzsolver}. We add three things: a precise
operational semantics for the engine, an LLM evaluation harness, and a precomputed
\emph{dead-state cache} that lets the harness distinguish provably unwinnable states from
mere losses.

\subsection{Predictive power of agentic benchmarks}
\label{sec:predictive}

Static benchmarks have repeatedly been saturated by scaling pretraining and reasoning
training data. ARC-AGI-1 \cite{arcagi1onmeasure} and ARC-AGI-2 \cite{arcagi2} resisted
direct memorization for years before frontier reasoning models began to make material
progress. ARC-AGI-3 \cite{arcagi3} pivots to interactive, turn-based environments where
information must be actively obtained.

HLG sits in the same family as ARC-AGI-3, but adds two mechanisms not exercised by ARC's
environments: \emph{persistent irreversible side-effects} (closing a required bridge) and
\emph{exact-state dead-end detection} via backward reachability analysis on the
state graph.

\subsection{Known limits of LRM fluid intelligence}
\label{sec:limits}

Modern Large Reasoning Models (LRMs) appear to acquire strong domain-specific reasoning
without acquiring corresponding general planning ability. They do well in domains where
training data densely samples the task space and where verifiable correctness signals
exist; they do poorly in domains where the relevant abstraction (here, backward
reachability on a state graph the model has never seen) cannot be retrieved from
training data.

HLG is designed so that surface-level imitation of past Bloxorz solutions does not
suffice: levels are non-trivial to enumerate publicly (the dead-state cache is computed
from the engine, not from any web corpus), and the dead-state vs.\ loss distinction
requires explicit reasoning that does not appear in standard puzzle-game text.

\subsection{Overfitting and memorization shortcuts}
\label{sec:overfit}

Like ARC-AGI-3, HLG mitigates memorization with a public/private split. Tutorial level
0 and the nine public levels (1--9) are intended as a demonstration set and are not
used for official scoring. The semi-private set (levels 10--24) is used for
public-API frontier-model leaderboards. The fully-private set (levels 25--33) is reserved
for competition use. Future revisions of HLG will expand the private set with newly
authored Bloxorz topologies that are guaranteed not to appear in any pre-2026 web
corpus.
